\chapter[]{Transport Layer Security (TLS)}

\section{Introduction to TLS}
This is a very famous and important security protocol. It has been invented more or less in the 90's, when internet was starting to grow to a large scale. More and more nodes were getting connected to internet, with users navigating to web servers.

Normally, we use an infrastructure that is composed of 5 layers to connect to the internet. There is a whole infrastructure that is in charge of transmitting the packets, where the \textbf{routers} are the routing elements. Not only the client, but also the server implements the 5 stacks.

At a certain point in time, the Netscape company anted to bring \textbf{secure payments} to internet, assuming that in the middle there could be \textbf{attackers}: on the communication lines, in the networking devices, anywhere. This is how TLS was born.

The \textbf{TLS is between the application layer and the transport layer}.

\section{TLS operation phases}
Lez mar 4 nov 2025\\
TLS protocol operates between 2 peers to secure the exchanged data. It works in phases.

The first phase is the handshake: there the security parameters are negotiated (authentication of the peers, in particular of the server, happens. for the server its mandatory, for the client its optional).

TLS phases:
\begin{itemize}
    \item TCP connection (3-way handshake)
    \item TLS handshake: authenticate server and optionally the client. Also negotiate the cryptographic algorithms to encrypt the data, and also protect the data for integrity and authentication (using a MAC and encryption). Keys are established
    \item (application) data transfer
    \item TLS teardown (closure): either invoked by application closing connection (normal closure) or due to error (during TLS handshake phase or the data transfer phase)
\end{itemize}

\paragraph{TLS handshake phase}
Random numbers are exchanged between client and server to generate keys. These numbers are called \textbf{client random} and \textbf{server random}. They get generated by scratch every handshake.

Client and server agree on a set of algorithms for protecting messages and key exchange; they are called \textbf{cipher suites}.

Key exchange/agreement means that client and server agree on a secret (called \textbf{master secret}), which is derived from another one, called \textbf{pre-master secret}. The master secret is basically a key, that is combined with client/server random to generate symmetric keys that protects the data. We have different keys for MAC and for encryption. Also different for the client and server.

\section{TLS architecture}
\section{TLS security services}
\section{TLS 1.2 server authentication only: steps and messages}

\section{Forward Secrecy in TLS}
An important reason why RSA is not used anymore in TLS is because it lacks the perfect forward secrecy.


----------------------------------------------------------------------
----------------------------------------------------------------------

\clearpage
